\documentclass[a4paper,12pt]{article}
\usepackage[utf8]{inputenc}
\usepackage[T1]{fontenc}
\usepackage[ngerman]{babel}
\usepackage{amsmath}
\usepackage{amssymb}
\usepackage{enumitem}
\usepackage{geometry}
\geometry{a4paper, margin=2.5cm}

\title{Einsendeaufgaben -- Lektion 1\\[0.5em]
\large Modul 61111: Mathematische Grundlagen}
\author{}
\date{}

\begin{document}
\maketitle

\section*{Aufgabe 1.3}

\begin{enumerate}[label=\alph*)]

\item $f: \mathbb{Z} \to \mathbb{Z} \times \mathbb{Z},\; x \mapsto (|x|, x^2)$

$f$ ist \textbf{nicht injektiv}, da z.\,B.\ $f(1) = (1,1) = f(-1)$ gilt.

\item $f: \mathbb{Z} \to \mathbb{Z} \times \mathbb{Z},\; x \mapsto (x, x)$

Aus $f(x) = f(y)$ folgt $(x, x) = (y, y)$, also muss $x = y$ gelten.
Demnach ist $f$ \textbf{injektiv}.

\item $f: \mathbb{Z} \to \mathbb{Z} \times \mathbb{Z},\; x \mapsto (x, 1)$

$f$ ist \textbf{nicht surjektiv}, da z.\,B.\ $(0, 0) \notin \mathrm{Bild}(f)$.

\item $f: \mathbb{Z} \to \mathbb{Z} \times \mathbb{Z},\; z \mapsto
\begin{cases}
  \left(\dfrac{z}{2},\, 0\right) & \text{wenn } 2 \mid z, \\[6pt]
  \left(0,\, \dfrac{z+1}{2}\right) & \text{wenn } 2 \nmid z.
\end{cases}$

\medskip
\textbf{Behauptung:}
$\mathrm{Bild}(f) = \{(z, 0) \mid z \in \mathbb{Z}\} \cup \{(0, z) \mid z \in \mathbb{Z}\}$.

\medskip
\textit{Richtung} ``$\subseteq$'':

Sei $z \in \mathbb{Z}$ beliebig.
\begin{itemize}
  \item Wenn $2 \mid z$ gilt, gibt es ein $x \in \mathbb{Z}$ mit $2x = z$.
        Dann ist $f(z) = \left(\frac{z}{2}, 0\right) \in \{(x, 0) \mid x \in \mathbb{Z}\}$.
  \item Wenn $2 \nmid z$ gilt, gibt es ein $x \in \mathbb{Z}$ mit $2x = z + 1$.
        Dann ist $f(z) = \left(0, \frac{z+1}{2}\right) \in \{(0, x) \mid x \in \mathbb{Z}\}$.
\end{itemize}

\textit{Richtung} ``$\supseteq$'':

Sei $z \in \mathbb{Z}$ beliebig.
\begin{itemize}
  \item $(z, 0) = \left(\frac{2z}{2}, 0\right) = f(2z) \in \mathrm{Bild}(f)$,
        da $2 \mid 2z$.
  \item $(0, z) = \left(0, \frac{(2z-1)+1}{2}\right) = f(2z-1) \in \mathrm{Bild}(f)$,
        da $2 \nmid (2z-1)$.
\end{itemize}

Aus ``$\subseteq$'' und ``$\supseteq$'' folgt
$\mathrm{Bild}(f) = \{(z, 0) \mid z \in \mathbb{Z}\} \cup \{(0, z) \mid z \in \mathbb{Z}\}$. \qed

\end{enumerate}

\end{document}
